% ------------------------------------------------------------
% Chapter 3: Geodesic Equations and Test Particles
% ------------------------------------------------------------

\chapter{Geodesic Equations and Test Particles}
\label{chap:geodesics}

\section{Overview}

In standard general relativity, freely falling test particles follow
geodesics of the spacetime metric $g_{\mu\nu}$. In \rsvpabbr{}, the
presence of lamphrodynamic entropy flux modifies the connection and thus
the notion of free fall. The purpose of this chapter is to derive the
modified geodesic equation, analyze its consequences for classical
tests of general relativity (perihelion precession, light bending),
and discuss gravitational time dilation in lamphrodynamic backgrounds.

Throughout, we assume the modified Einstein equations from
\cref{chap:gravity-entropy}:
\[
  G_{\mu\nu} + \Lambda g_{\mu\nu} + \Theta_{\mu\nu} = 8\pi G\,T_{\mu\nu},
\]
where $\Theta_{\mu\nu}$ encodes lamphrodynamic contributions from
$(\PhiField,\vField,\SField)$ as derived in \cref{sec:theta-explicit}.

\section{Modified Connection and Effective Metric}

\subsection{Lamphrodynamic Correction}

We model the lamphrodynamic influence on geodesic motion by defining an
effective connection
\[
  \tilde{\Gamma}^{\lambda}_{\mu\nu}
  = \Gamma^{\lambda}_{\mu\nu}
  + \Delta^{\lambda}_{\mu\nu},
\]
where $\Gamma$ is the Levi--Civita connection of $g$ and
$\Delta^{\lambda}_{\mu\nu}$ is a small correction depending on
$\Theta_{\mu\nu}$ and derivatives of $\PhiField$, $\vField$, and $\SField$.
For definiteness, we adopt the model ansatz
\begin{equation}
\label{eq:delta-connection}
  \Delta^{\lambda}_{\mu\nu}
  := \alpha\,g^{\lambda\rho}
  \nabla_{(\mu}\bigl(\SField\,\vField_{\nu)}\bigr),
\end{equation}
where $\alpha$ is a coupling constant with dimensions of length, and
$(\cdot)_{(\mu\nu)}$ denotes symmetrization. More sophisticated forms of
$\Delta$ could be derived from the AKSZ action, but \eqref{eq:delta-connection}
captures the leading lamphrodynamic contribution in weak fields.

\subsection{Effective Geodesic Equation}

The modified geodesic equation becomes
\begin{equation}
\label{eq:modified-geodesic}
  \frac{d^2 x^\lambda}{d\tau^2}
  + \Gamma^\lambda_{\mu\nu}
  \frac{dx^\mu}{d\tau}\frac{dx^\nu}{d\tau}
  + \Delta^\lambda_{\mu\nu}
  \frac{dx^\mu}{d\tau}\frac{dx^\nu}{d\tau}
  = 0.
\end{equation}
When $\Delta\equiv 0$, we recover the usual geodesic equation.

\begin{remark}[Interpretation]
The correction term encodes the effect of entropy flow and lamphrodynamic
gradients on inertial motion. In heuristic terms, test particles
``feel'' an additional force resulting from entropy flux, analogous to
thermodynamic drag or buoyancy effects in continuum media, but in a fully
covariant setting.
\end{remark}

\section{Modified Raychaudhuri Equation}

\subsection{Classical Form}

For a congruence of timelike geodesics with tangent vector $u^\mu$, the
classical Raychaudhuri equation is
\[
  \frac{d\theta}{d\tau}
  = -\frac{1}{3}\theta^2 - \sigma_{\mu\nu}\sigma^{\mu\nu}
  + \omega_{\mu\nu}\omega^{\mu\nu}
  - R_{\mu\nu} u^\mu u^\nu,
\]
where $\theta$ is expansion, $\sigma$ shear, and $\omega$ vorticity.

\subsection{Lamphrodynamic Correction}

In \rsvpabbr{}, the Ricci tensor $R_{\mu\nu}$ is replaced by an effective
Ricci tensor $\tilde{R}_{\mu\nu}$ defined by
\[
  \tilde{R}_{\mu\nu}
  := R_{\mu\nu} + \Theta_{\mu\nu}.
\]
Thus, the modified Raychaudhuri equation becomes
\begin{equation}
\label{eq:raychaudhuri-modified}
  \frac{d\theta}{d\tau}
  = -\frac{1}{3}\theta^2 - \sigma_{\mu\nu}\sigma^{\mu\nu}
  + \omega_{\mu\nu}\omega^{\mu\nu}
  - (R_{\mu\nu} + \Theta_{\mu\nu}) u^\mu u^\nu.
\end{equation}
The extra term $-\Theta_{\mu\nu}u^\mu u^\nu$ may either enhance or reduce
focusing of geodesics depending on the sign and magnitude of entropy
gradients and lamphrodynamic flux.

\section{Perihelion Precession}

\subsection{Standard GR Result}

In general relativity, the perihelion of Mercury advances by approximately
$43''$ per century due to relativistic corrections. Observational agreement
is excellent.

\subsection{Lamphrodynamic Correction}

To leading order in the weak--field, slow--motion limit, the correction to
the effective potential due to $\Delta^\lambda_{\mu\nu}$ yields a modified
perihelion advance of the form
\[
  \Delta\phi_{\mathrm{RSVP}}
  \approx \frac{6\pi G M}{a(1-e^2)c^2}
  \bigl(1 + \epsilon_{\mathrm{RSVP}}\bigr),
\]
where $a$ and $e$ are the semi--major axis and eccentricity, and
$\epsilon_{\mathrm{RSVP}}$ is a small dimensionless parameter proportional
to $\alpha\,\nabla \SField \cdot \vField$ evaluated along the orbit.

Empirical bounds on $\epsilon_{\mathrm{RSVP}}$ arise from detailed fits
to planetary ephemerides; current observational precision suggests
$\epsilon_{\mathrm{RSVP}} \lesssim 10^{-5}$, depending on the specific
model of $\SField$ and $\vField$.

\section{Light Bending}

\subsection{Null Geodesics}

For null geodesics ($ds^2=0$), the same modified geodesic equation
\eqref{eq:modified-geodesic} applies. To leading order, the light--bending
angle near a massive body is
\[
  \delta\phi_{\mathrm{RSVP}}
  \approx \frac{4GM}{c^2 b}
  \bigl(1 + \epsilon_{\mathrm{RSVP}}\bigr),
\]
where $b$ is the impact parameter. Current measurements of gravitational
lensing by the Sun impose tight constraints on $\epsilon_{\mathrm{RSVP}}$,
potentially at the level of $10^{-4}$ or better.

\section{Gravitational Time Dilation}

Lamphrodynamic corrections modify redshift and time dilation through
$\Delta^\lambda_{\mu\nu}$ and $\Theta_{\mu\nu}$. The proper time ratio
between two stationary observers at different potentials receives a
multiplicative correction of the form
\[
  \frac{d\tau_1}{d\tau_2}
  = \sqrt{\frac{g_{00}(x_1)}{g_{00}(x_2)}}
  \bigl(1 + \epsilon_{\mathrm{RSVP}}\bigr),
\]
where again $\epsilon_{\mathrm{RSVP}}$ is controlled by entropy gradients
and lamphrodynamic fluxes. Existing experimental tests of gravitational
redshift are sufficiently precise that bounds on $\epsilon_{\mathrm{RSVP}}$
are likely stringent, though a systematic analysis awaits detailed
phenomenological modeling.

\section{Summary}

Modifications to geodesic motion provide one of the clearest empirical
handles on lamphrodynamic gravity. Classical tests of relativity --- 
perihelion precession, light bending, gravitational time dilation ---
place strong constraints on the magnitude of lamphrodynamic corrections.
Nonetheless, small deviations remain possible and could be detected by
higher--precision experiments, providing a direct test of the \rsvpabbr{}
framework.

In the next chapter (\cref{chap:friedmann-like}), we turn to cosmological
implications, deriving Friedmann--like equations from homogeneous and
isotropic lamphrodynamic configurations and exploring the entropic
contribution to cosmic expansion.

\clearpage

